% Part: incompleteness
% Chapter: arithmetization-syntax
% Section: proofs-in-lk

\documentclass[../../../include/open-logic-section]{subfiles}

\begin{document}

\olfileid{inc}{art}{plk}
\olsection{$\Log{LK}$ وچ \usetoken{P}{derivation}}

\begin{explain}
!!{derivation}s نوں عددی روپ دین لئی سانوں !!{derivation}s نوں
عدداں نال ظاہر کرنا پینا اے۔ !!{derivation}s سیکوئنٹاں دے درخت نیں
تے ہر استنباط دا اک لیبل وی اے۔ ایس لئی بازگشتی طریقہ سِدھا اے۔
اسیں ہر !!{derivation} نوں اک ٹپل نال ظاہر کردے آں، جیہدے جز
آخرلا سیکوئنٹ، استنباط دا لیبل تے آخری استنباط دے مقدمیاں تک
پہنچن والے ذیلی !!{derivation}s دے عددی روپ نیں۔
\end{explain}

\begin{defn}
جے $\Gamma$، !!{sentence}s دا متناہی سلسلہ اے، $\Gamma =
\tuple{!A_1, \dots, !A_n}$، تاں
$\Gn{\Gamma} = \tuple{\Gn{!A_1},
  \dots, \Gn{!A_n}}$ اے۔

جے $\Gamma \Sequent \Delta$
اک سیکوئنٹ اے، تاں
$\Gamma \Sequent \Delta$ دا
گوڈل نمبر ایہ اے:
\[
\Gn{\Gamma \Sequent \Delta} = \tuple{\Gn{\Gamma}, \Gn{\Delta}}
\]

جے $\pi$، $\Log{LK}$ وچ
!!a{derivation} اے، تاں
$\Gn{\pi}$ دی تعریف ایویں کردے آں:
\begin{enumerate}
\item جے $\pi$ صرف ابتدائی
  سیکوئنٹ $\Gamma \Sequent \Delta$
  اُتے مشتمل ہووے، تاں
  $\Gn{\pi}$ ایہ اے:
  \[
    \tuple{0, \Gn{\Gamma \Sequent \Delta}}.
  \]
\item جے $\pi$ دا آخرلا استنباط
  اک یا دو مقدمیاں والا ہووے،
  اوہدا نتیجہ
  $\Gamma \Sequent \Delta$
  ہووے، تے $\pi_1$ تے
  $\pi_2$ اوہ براہِ راست ذیلی ثبوت ہون جیہڑے آخری استنباط دے
  مقدمیاں اُتے ختم ہوندے نیں، تاں $\Gn{\pi}$ دی قدر بالترتیب ایہ اے:
  \begin{align*}
    & \tuple{1, \Gn{\pi_1}, \Gn{\Gamma \Sequent \Delta}, k} \text{ یا}\\
    & \tuple{2, \Gn{\pi_1}, \Gn{\pi_2}, \Gn{\Gamma \Sequent \Delta}, k},
  \end{align*}
  ایتھے $k$ دی قدر، آخری استنباط وچ ورتے قاعدے مطابق،
  ہِٹھلے جدول توں ملدی اے:

  \begin{tabular}{lccccccc}
    \text{قاعدہ:} & \LeftR{\Weakening} & \RightR{\Weakening} &
    \LeftR{\Contraction} & \RightR{\Contraction} &
    \LeftR{\Exchange} & \RightR{\Exchange} \\
    $k$: & 1 & 2 & 3 & 4 & 5 & 6 \\[2ex]
    \text{قاعدہ:} & \LeftR{\lnot} & \RightR{\lnot} &
    \LeftR{\land} & \RightR{\land} &
    \LeftR{\lor} & \RightR{\lor} \\
    $k$: & 7 & 8 & 9 & 10 & 11 & 12 \\[2ex]
    \text{قاعدہ:} & \LeftR{\lif} & \RightR{\lif} &
    \LeftR{\lforall} & \RightR{\lforall} &
    \LeftR{\lexists} & \RightR{\lexists} \\
    $k$: & 13 & 14 & 15 & 16 & 17 & 18 \\[2ex]
    \text{قاعدہ:} & \Cut & = \\
    $k$: & 19 & 20
  \end{tabular}
\end{enumerate}
\end{defn}

\begin{ex}
ایہ سوکھا جیہا !!{derivation} ویکھو:
\begin{prooftree}
  \Axiom$!A \fCenter !A$
  \RightLabel{\LeftR{\land}}
  \UnaryInf$!A \land !B \fCenter !A$
  \RightLabel{\RightR{\lif}}
  \UnaryInf$\fCenter (!A \land !B) \lif !A$
\end{prooftree}
ابتدائی سیکوئنٹ دا گوڈل نمبر
$p_0 = \tuple{0,
  \Gn{!A \Sequent !A}}$
ہوئے گا۔ $\LeftR{\land}$ دے
نتیجے اُتے ختم ہون والے
!!{derivation} دا گوڈل نمبر
$p_1 =
    \tuple{1, p_0, \Gn{!A \land !B \Sequent !A}, 9}$
ہووے گا ($1$ ایس لئی کہ
$\LeftR{\land}$ دا اک مقدمہ اے؛
نتیجے~$!A \land !B \Sequent !A$
دا گوڈل نمبر؛ تے $9$،
$\LeftR{\land}$ نوں کوڈ کرن
والا نمبر اے)۔ فیر پورے
!!{derivation} دا گوڈل نمبر
$\tuple{1, p_1, \Gn{\Sequent (!A \land !B) \lif !A)}, 14}$
اے؛ یعنی
\[
\tuple{1, \tuple{1, \tuple{0, \Gn{!A \Sequent !A)}}, \Gn{!A \land !B \Sequent !A}, 9},
  \Gn{\Sequent (!A \land !B) \lif !A}, 14}.
\]
\end{ex}

\begin{explain}
!!{derivation}s دی عددی روپ طے کر کے سانوں ایہ وی دکھانا
پینا اے کہ اسیں ایہناں
!!{derivation}s اُتے بنیادی بازگشتی
ڈھنگ نال عمل کر سکدے آں تے انہاں
دیاں ضروری خاصیتاں تے تعلقاں
نوں وی ایس طرح بیان کر سکدے آں۔
کجھ عمل سادے نیں: مثال لئی
!!a{derivation} دے گوڈل
نمبر~$p$ توں
$\fn{EndSeq}(p) = (p)_{(p)_0+1}$
اوس دے آخرلے سیکوئنٹ دا
گوڈل نمبر دیندا اے تے
$\fn{LastRule}(p) = (p)_{(p)_0+2}$
اوس دے آخری قاعدے دا
کوڈ دیندا اے۔ خاصیت
$\fn{Sequent}(s)$، جیہدی
تعریف ایہ اے،
\[
  \len{s} = 2 \land \bforall{i<\len{(s)_0} + \len{(s)_1}}{\fn{Sent}(((s)_0 \concat (s)_1)_i)}
\]
$s$ لئی عین اوس ویلے
سچ اے جدوں $s$،
!!{sentence}s توں بنے
سیکوئنٹ دا گوڈل نمبر
ہووے۔ کجھ ہور عمل کِتے ودھ اوکھے نیں۔ ایتھے اسیں گھٹ توں گھٹ اوہناں دا خاکہ دساں گے۔
مقصد ایہ دکھانا اے کہ
«$\pi$، $!A$ دا
!!a{derivation} $\Gamma$
توں اے» والا تعلق، $\pi$ تے
$!A$ دے گوڈل نمبراں
دا بنیادی بازگشتی تعلق اے۔
\end{explain}

\begin{prop}\ollabel{prop:followsby}
خاصیت $\fn{Correct}(p)$ اُدوں تے صرف اُدوں قائم اے جدوں
!!{derivation}~$\pi$ دا آخری استنباط درست ہووے، جتھے اوس اخذ دا
گوڈل نمبر~$p$ اے۔ ایہ خاصیت بنیادی بازگشتی اے۔
\end{prop}

\begin{proof}
$\Gamma \Sequent \Delta$
ابتدائی سیکوئنٹ اے جے یا تاں
کوئی !!a{sentence}~$!A$
ایہوجیہا ہووے کہ
$\Gamma \Sequent \Delta$،
$!A \Sequent !A$ ہووے؛
یا کوئی ٹرم~$t$ ایہوجیہا
ہووے کہ
$\Gamma \Sequent \Delta$،
$\emptyset \Sequent \eq[t][t]$
ہووے۔ گوڈل نمبراں دے لحاظ نال
$\fn{InitSeq}(s)$ عین
اوس ویلے سچ اے جدوں
\begin{align*}
\bexists{x < s}{} (\fn{Sent}(x) & \land
s = \tuple{\tuple{x},\tuple{x}})
\lor {}\\
\bexists{t<s}{} (\fn{Term}(t) & \land
s = \tuple{0, \tuple{\Gn{{\eq}(} \concat t \concat \Gn{,} \concat t \concat \Gn{)}}}).
\end{align*}

سانوں ایہ وی دکھانا
پینا اے کہ استنباط دے
ہر قاعدے~$R$ لئی
تعلق $\fn{FollowsBy}_R(p)$
بنیادی بازگشتی اے۔
ایتھے
$\fn{FollowsBy}_R(p)$
عین اوس ویلے سچ اے
جدوں $p$، کسے
!!{derivation}~$\pi$
دا گوڈل نمبر ہووے،
تے~$\pi$ دا آخرلا
سیکوئنٹ، قاعدہ~$R$ ٹھیک طرح لاؤن نال~$\pi$ دے براہِ راست ذیلی
!!{derivation}s دے نتیجیاں توں نکلدا ہووے۔

اک سوکھی مثال \RightR{\land} دا قاعدہ اے۔ جے $\pi$ اک درست
$\RightR{\land}$ استنباط اُتے ختم ہووے، تاں اوہ ایویں دِسے گا:
\begin{prooftree}
  \AxiomC{}
  \RightLabel{$\pi_1$}
  \Deduce$\Gamma \fCenter \Delta, !A$

    \AxiomC{}
    \RightLabel{$\pi_2$}
    \Deduce$\Gamma \fCenter \Delta, !B$

  \RightLabel{\RightR\land}
  \BinaryInf$\Gamma \fCenter \Delta, !A \land !B$
\end{prooftree}
ایس لئی
!!{derivation} $\pi$
وچ آخری استنباط
$\RightR{\land}$ دا
درست اطلاق اُدوں تے صرف اُدوں اے جدوں
!!{sentence}s دے
سلسلے $\Gamma$
تے $\Delta$، نالے
دو !!{sentence}s
$!A$ تے~$!B$
ایہوجیہے ہون کہ
$\pi_1$ دا آخرلا
سیکوئنٹ
$\Gamma \Sequent \Delta, !A$
ہووے، $\pi_2$ دا
آخرلا سیکوئنٹ
$\Gamma \Sequent \Delta, !B$
ہووے، تے $\pi$
دا آخرلا سیکوئنٹ
$\Gamma \Sequent \Delta, !A \land !B$
ہووے۔ ہن بس ایس نوں
گوڈل نمبراں وچ
بیان کرنا اے۔ جے
$s = \Gn{\Gamma \Sequent \Delta}$
ہووے، تاں
$(s)_0 = \Gn{\Gamma}$
تے $(s)_1 = \Gn{\Delta}$
ہوندے نیں۔ ایس لئی
$\fn{FollowsBy}_{\RightR{\land}}(p)$
عین اوس ویلے سچ اے
جدوں
\begin{align*}
& \bexists{g < p}{\bexists{d < p}{\bexists{a < p}{\bexists{b < p}{\quad}}}} \\
& \qquad \fn{EndSequent}(p) =
  \tuple{g, d \concat
    \tuple{\Gn{(} \concat a \concat \Gn{\land} \concat b \concat \Gn{)}}} \land {} \\
& \qquad \fn{EndSequent}((p)_1) = \tuple{g, d \concat \tuple{a}} \land {}\\
& \qquad \fn{EndSequent}((p)_2) = \tuple{g, d \concat \tuple{b}} \land {}\\
& \qquad (p)_0 = 2 \land \fn{LastRule}(p) = 10.
\end{align*}
ہر سطر دی گل بالترتیب ایہ اے کہ
«سلسلہ~($\Gamma$)،
جیہدا گوڈل نمبر~$g$ اے،
سلسلہ~($\Delta$)،
جیہدا گوڈل نمبر~$d$ اے،
!!a{formula}~($!A$)،
جیہدا گوڈل نمبر~$a$ اے،
تے !!a{formula}~($!B$)،
جیہدا گوڈل نمبر~$b$ اے،
موجود نیں»، تے
نال ایہ کہ
«$\pi$ دا آخرلا
سیکوئنٹ
$\Gamma \Sequent \Delta, !A
\land !B$ اے»،
«$\pi_1$ دا آخرلا
سیکوئنٹ
$\Gamma \Sequent \Delta,
!A$ اے»،
«$\pi_2$ دا آخرلا
سیکوئنٹ
$\Gamma \Sequent \Delta, !B$
اے»، تے
«$\pi$ دیاں دو
براہِ راست ذیلی
اخذاں نیں تے آخری
استنباطی قاعدہ
$\RightR\land$
اے (جیہدا نمبر~$10$
اے)»۔

$\pi$ وچ آخری استنباط
$\RightR\lexists$
دا درست اطلاق اُدوں تے صرف اُدوں اے جدوں
سلسلے $\Gamma$
تے $\Delta$، اک
!!a{formula}~$!A$،
اک متغیر~$x$،
تے اک ٹرم~$t$
ایہوجیہے ہون کہ
$\pi$ دا آخرلا
سیکوئنٹ
$\Gamma \Sequent \Delta, \lexists[x][!A]$
تے $\pi_1$ دا
آخرلا سیکوئنٹ
$\Gamma \Sequent \Delta,
\Subst{!A}{t}{x}$
ہووے۔ ایس لئی گوڈل
نمبراں دے لحاظ نال
$\fn{FollowsBy}_{\RightR\lexists}(p)$
عین اوس ویلے سچ اے
جدوں
\begin{align*}
  & \bexists{g<p}{\bexists{d<p}{\bexists{a<p}{\bexists{x<p}{\bexists{t<p}{\quad}}}}}\\
  & \qquad \fn{EndSequent}(p) =
  \tuple{
    g,
    d \concat \tuple{\Gn{\lexists} \concat x \concat a}
  } \land {}\\
  & \qquad \fn{EndSequent}((p)_1) =
  \tuple{
    g,
    d \concat \tuple{\fn{Subst}(a, t, x)}
  } \land {}\\
  & \qquad (p)_0 = 1 \land \fn{LastRule}(p) = 18.
\end{align*}
فیر اسیں $\fn{Correct}(p)$ دی تعریف ایویں کردے آں:
\begin{multline*}
  \fn{Sequent}(\fn{EndSequent}(p)) \land {}\\
  [(\fn{LastRule}(p) = 1 \land
  \fn{FollowsBy}_{\LeftR\Weakening}(p)) \lor \dots \lor {}\\
  (\fn{LastRule}(p) = 20 \land \fn{FollowsBy}_{\eq}(p)) \lor {}\\
  (p)_0 = 0 \land \fn{InitialSeq}(\fn{EndSequent}(p))]
\end{multline*}
پہلی سطر یقینی بناندی
اے کہ~$d$ دا آخرلا
سیکوئنٹ واقعی
!!{sentence}s دا
سیکوئنٹ اے۔ آخری
سطر اوس صورت لئی اے
جدوں $p$ بس اک
ابتدائی سیکوئنٹ ہووے۔
\end{proof}

\begin{prob}
ہِٹھلیاں خاصیتاں
\olref[inc][art][plk]{prop:followsby}
دی طرز اُتے تعریف کرو:
\begin{enumerate}
\item $\fn{FollowsBy}_{\Cut}(p)$،
\item $\fn{FollowsBy}_{\LeftR\lif}(p)$،
\item $\fn{FollowsBy}_{\eq}(p)$،
\item $\fn{FollowsBy}_{\RightR{\lforall}}(p)$۔
\end{enumerate}
آخری صورت لئی ایہ وی وکھاؤ کہ
گوڈل نمبر~$p$
والے !!{derivation}
دا آخری استنباط
ایگن متغیر والی
شرط پوری کردا
اے یا نہیں، ایہ
بنیادی بازگشتی
ڈھنگ نال پرکھیا
جا سکدا اے۔ یعنی
ایگن متغیر~$a$،
جیہڑا $\RightR{\lforall}$
استنباط دا اے،
آخرلے سیکوئنٹ
وچ نہ آوے۔
\end{prob}

\begin{prop}
\ollabel{prop:deriv}
تعلق $\fn{Deriv}(p)$ بنیادی بازگشتی اے۔ ایہ قائم اے جے $p$
کسے درست !!{derivation}~$\pi$ دا گوڈل نمبر ہووے۔
\end{prop}

\begin{proof}
!!^a{derivation}~$\pi$
اوس ویلے درست اے
جدوں اوہدا ہر استنباط کسے قاعدے
دا درست اطلاق
ہووے؛ یعنی اوہدے ذیلی !!{derivation}s
وچوں ہر اک
درست استنباط اُتے
ختم ہووے۔ ایس لئی
$\fn{Deriv}(d)$
عین اوس ویلے سچ
اے جدوں
\[
\bforall{i<\len{\fn{SubtreeSeq}(p)}}{\fn{Correct}((\fn{SubtreeSeq}(p))_i}.
\]
\end{proof}

\begin{prop}
فرض کرو $\Gamma$،
!!{sentence}s دا
بنیادی بازگشتی
سیٹ اے۔ تاں تعلق
$\Prf[\Gamma](x, y)$،
جیہڑا ایہ بیان
کردا اے کہ
«$x$،
!!a{derivation}~$\pi$
دا کوڈ اے جیہدا
آخرلا سیکوئنٹ
$\Gamma_0 \Sequent !A$
اے، جتھے متناہی $\Gamma_0 \subseteq \Gamma$
ہووے، تے
$y$،~$!A$ دا
گوڈل نمبر اے»،
بنیادی بازگشتی
اے۔
\end{prop}

\begin{proof}
فرض کرو
«$y \in \Gamma$»
بنیادی بازگشتی
محمول~$R_\Gamma(y)$
نال دتا گیا اے۔
سانوں ایہ دکھانا
اے کہ
$\Prf[\Gamma](x, y)$،
جیہڑا عین اوس
ویلے سچ اے جدوں
$y$ کسے
جملے~$!A$ دا
گوڈل نمبر ہووے
تے $x$، $\Log{LK}$ وچلے کسے !!{derivation} دا کوڈ ہووے جیہدا
آخرلا سیکوئنٹ
$\Gamma_0 \Sequent !A$
ہووے،
بنیادی بازگشتی
اے۔

پچھلے قضیے مطابق خاصیت
$\fn{Deriv}(x)$،
جیہڑی عین اوس
ویلے سچ اے جدوں
$x$ کسے درست
!!{derivation}~$\pi$
دا کوڈ ہووے
جیہڑا $\Log{LK}$
وچ اے،
بنیادی بازگشتی
اے۔ جے $x$
ایہوجیہا کوڈ
اے، تاں
$\fn{EndSequent}(x)$،
~$\pi$ دے
آخرلے سیکوئنٹ
دا کوڈ اے، تے
ایس لئی
$(\fn{EndSequent}(x))_0$
آخرلے سیکوئنٹ
دے کھبے پاسے
دا کوڈ، تے
$(\fn{EndSequent}(x))_1$
سجے پاسے دا
کوڈ اے۔ ایس
لئی اسیں
«$\pi$ دے
آخرلے سیکوئنٹ
دا سجا پاسا
صرف~$!A$ اے»
نوں
$\len{(\fn{EndSequent}(x))_1} = 1 \land ((\fn{EndSequent}(x))_1)_0 =
x$
نال بیان کر
سکدے آں۔
$\pi$ دے
آخرلے سیکوئنٹ
دا کھبا پاسا آپے متناہی اے۔ سانوں صرف ایہ بیان کرنا اے کہ
اوہدا ہر جملہ~$\Gamma$ وچ اے۔ ایس لئی
$\Prf[\Gamma](x, y)$
نوں ایس طرح
تعریف کر سکدے
آں:
\begin{align*}
\Prf[\Gamma](x, y) \defiff {}&
\fn{Deriv}(x) \land {} \\
& \bforall{i <
  \len{(\fn{EndSequent}(x))_0}}{R_\Gamma(((\fn{EndSequent}(x))_0)_i)} \land {}\\
& \len{(\fn{EndSequent}(x))_1} = 1 \land ((\fn{EndSequent}(x))_1)_0 = y.
\end{align*}
\end{proof}

\end{document}
